% 固定模型之上的智能体团队：The One 两个月框架实验（中文版）
% 编译：tectonic -X compile main.tex（XeLaTeX；BibTeX 自动运行）
% 字体：fontset=mac 使用 macOS 自带中文字体；在 Linux 或 Windows 上改为 fontset=fandol（TeX Live 自带）。
\documentclass[UTF8,zihao=5,fontset=mac]{ctexart}

\usepackage[a4paper,textwidth=152mm,textheight=232mm,headsep=8mm,footskip=11mm]{geometry}
\usepackage{amsmath}
\usepackage{unicode-math}
\setmainfont{texgyretermes}[Extension=.otf,UprightFont=*-regular,BoldFont=*-bold,
  ItalicFont=*-italic,BoldItalicFont=*-bolditalic]
\setmathfont{texgyretermes-math.otf}
\setsansfont{texgyreheros}[Extension=.otf,UprightFont=*-regular,BoldFont=*-bold,
  ItalicFont=*-italic,BoldItalicFont=*-bolditalic,Scale=0.9]
\setmonofont{Inconsolatazi4-Regular.otf}[BoldFont=Inconsolatazi4-Bold.otf,Scale=0.9]

\usepackage{booktabs,tabularx,longtable,array,makecell}
\usepackage[table]{xcolor}
\usepackage{graphicx}
\usepackage[font=small,labelfont=bf,labelsep=quad,skip=5pt]{caption}
\usepackage{enumitem}
\setlist{itemsep=1pt,topsep=3pt,parsep=1pt}
\usepackage{tikz}
\usetikzlibrary{arrows.meta,positioning,fit,calc,backgrounds,decorations.pathreplacing}
\usepackage{pgfplots}
\pgfplotsset{compat=1.18}
\usepackage[numbers,sort&compress]{natbib}
\setlength{\bibsep}{2pt plus 1pt}
\usepackage{fancyhdr}
\usepackage[hidelinks]{hyperref}

\ctexset{
  section={format=\large\bfseries\raggedright, beforeskip={12pt plus 3pt minus 2pt},
    afterskip={6pt plus 1pt}},
  subsection={format=\normalsize\bfseries\raggedright, beforeskip={9pt plus 2pt minus 2pt},
    afterskip={4pt plus 1pt}},
  abstractname={摘\quad 要},
  bibname={参考文献},
  appendixname={附录},
}
\newcommand{\para}[1]{\par\vspace{3pt plus 1pt}\noindent\textbf{#1}\quad}

\definecolor{ours}{HTML}{C2502E}
\definecolor{oursLight}{HTML}{E9A58F}
\definecolor{other}{HTML}{8A8F98}
\definecolor{ink}{HTML}{2B2F36}
\definecolor{boxfill}{HTML}{F4F1EC}
\definecolor{ruleGray}{HTML}{B9BEC6}

\newcolumntype{L}{>{\raggedright\arraybackslash}X}
\renewcommand{\arraystretch}{1.15}
\newcommand{\code}[1]{\texttt{#1}}

\pagestyle{fancy}
\fancyhf{}
\renewcommand{\headrulewidth}{0.3pt}
\fancyhead[L]{\small 固定模型之上的智能体团队}
\fancyhead[R]{\small 张子阳，2026}
\fancyfoot[C]{\small\thepage}
\fancypagestyle{plain}{\fancyhf{}\renewcommand{\headrulewidth}{0pt}\fancyfoot[C]{\small\thepage}}

\hypersetup{
  pdftitle={固定模型之上的智能体团队：The One 两个月框架实验},
  pdfauthor={张子阳 (Ziyang Zhang)},
  pdfsubject={在固定模型上对智能体框架与团队结构做的对照实验},
  pdfkeywords={大语言模型智能体, 编程智能体, 多智能体系统, 智能体框架, 上下文工程, 提示词缓存,
    Terminal-Bench, 评估方法}
}

\begin{document}

\thispagestyle{plain}
\begin{center}
  {\small 预印本，2026 年 11 月\quad\textperiodcentered\quad DOI：\href{https://doi.org/10.5281/zenodo.24410340}{10.5281/zenodo.24410340}}\\[10pt]
  {\LARGE\bfseries 固定模型之上的智能体团队：\\[4pt] The One 两个月框架实验\par}
  \vspace{12pt}
  {\large 张子阳}\\[3pt]
  {\normalsize The Zeroth}\\[1pt]
  {\small\url{https://the-zeroth.com}}\\[6pt]
  {\footnotesize 英文版：\textit{Teams on a Fixed Model: Two Months of Harness Experiments with The One}}
\end{center}
\vspace{4pt}

\begin{abstract}
\noindent
模型固定时，一个智能体团队框架（harness）还能带来多少提升？本文报告 2026 年 9 月 28 日至
11 月 20 日共八周在 The One 上做的对照实验。The One 是一款在个人电脑上运行语言模型智能体团队的
产品。三个实验项目先后进行：V1 在我们自建的任务集上进行（120 多批）；V8 通过 Harbor 框架在
Terminal-Bench~4.0 上进行（315 个作业、6,223 次已评分试验，其中 30 道题留出、从未读过）；最后按
Artificial Analysis 的协议在 Terminal-Bench~4.1 上跑了一次。最终保留的配置是单个领队智能体，用
GPT-6.2 Sol、最高推理强度（max）。它在能运行的 63 道题上得分 60.9\%，把 3 道需要 GPU 的题记为失败时得分
58.1\%，与模型服务商自家的智能体（Codex）在同一模型、同一强度下持平，同时输入 token 少约 30\%，
按公开标价每次尝试 1.35 美元。本文有三个主要发现。第一，只追加的上下文把成本降低三分之一到一半，
准确率不变。第二，在单个上下文装得下的任务上，一个能力强的负责者不输于我们搭建的任何团队结构，
而同一模型给出的第二意见只会确认共同的误读，抓不住它。第三，模型在决策点上的选择只有在事实摆到它
面前时才会改变，建议从未改变它。多数在设计任务上胜出的手段到了留出任务上就失效了；本文介绍了发现
这些问题的留出纪律，并完整报告了负面结果。

\vspace{6pt}\noindent
\textbf{关键词：}大语言模型智能体；编程智能体；多智能体系统；智能体框架；上下文工程；
提示词缓存；Terminal-Bench；评估方法
\end{abstract}

\section{引言}
\label{sec:intro}

编程智能体排行榜现在排的是“智能体 + 模型”的组合，而不只是模型本身
\citep{tbench2026board,aa2026coding,kapoor2026hal}，许多产品包装的也是同样几个前沿模型。对于做这层包装的人来说，
框架（harness）指的是模型外围的一切：提示词、工具、上下文管理、团队结构和运行时检查。实际要回答的问题是：
模型固定以后，框架还能带来什么。已有的答案并不一致。多智能体框架报告了角色扮演团队、辩论和投票带来的提升
\citep{wu2024autogen,hong2024metagpt,qian2024chatdev,du2024debate,li2024moreagents}；对这类系统的失败分析发现，
它们相对单个智能体的提升往往很小 \citep{cemri2025mast}；业界实践者有的主张共享上下文的单个智能体
\citep{yan2025dontbuild}，有的主张由领队把工作分给并行的助手 \citep{anthropic2025multiagent}。

我们开发 The One，这款产品在个人电脑或 Linux 服务器上运行语言模型智能体团队。2026 年 9 月 28 日至 11 月 20 日，
我们在它上面先后进行了三个实验项目，回答四个问题：
\begin{enumerate}[label=\textbf{RQ\arabic*},leftmargin=3.2em]
  \item 上下文管理在成本、延迟和准确率上能换来什么？
  \item 团队结构能否胜过一个能力强的单个智能体？在哪些任务上？
  \item 框架能否在运行失败的那些关键点上引导模型的判断？
  \item 由此得到的配置，与模型服务商自家的智能体相比处在什么位置？
\end{enumerate}

\para{贡献}本文的贡献主要是一份记录和一套方法，而不是一个系统。
\begin{itemize}[leftmargin=1.4em]
  \item \textbf{对照实验记录。}在自建任务集上做了 120 多批实验（V1）；在 Terminal-Bench~4.0 上跑了 315 个 Harbor
    作业、6,223 次已评分试验（V8）；按 Artificial Analysis（AA）的协议在 Terminal-Bench~4.1 上跑了一次。每批实验
    开跑之前，都先写下规则、分支和我们自己的预测。全部实验的索引见附录~\ref{app:index}。
  \item \textbf{一套分辨率已知的留出纪律。}30 道 Terminal-Bench 题目从未打开；每个实验组都和它的基线在同一时段
    一起运行；差异通过题内置换检验来判断；得零分的试验在空闲机器上重评。这套纪律抓出了我们自己的大部分
    “假胜利”，我们也给出了它能分辨的最小效应（第~\ref{sec:method}~节）。
  \item \textbf{实测效应，大多是负面的。}上下文手段对成本和延迟的影响（第~\ref{sec:context}~节）；团队结构和同模型
    第二意见共有的失败模式（第~\ref{sec:structure}~节）；以及一条引导模型的规律：提示改变过程，事实改变选择
    （第~\ref{sec:judgment}~节）。
  \item \textbf{公开的适配器。}Harbor 适配器和完整的运行命令已经公开 \citep{zhang2026adapter}；该配置在
    Terminal-Bench~4.1 上与服务商自家的智能体持平，输入 token 更少（第~\ref{sec:tb41}~节）。
\end{itemize}

\para{主要发现}只追加的提示词把每次运行的成本降低了三分之一到一半，得分不变；之后的上下文手段又降低了
7\% 到 20\%，但都没有改变准确率。在单个上下文装得下的任务上，一个能力强的负责者不输于我们搭建的任何团队，
成本通常只是团队的一小部分，原因是剩下的失败都是单一的误读，同一模型的每个副本都会犯同样的错。
给模型提建议的文字改变了它下一步做什么，却没有改变它的选择；把事实和测量结果摆到它面前，才改变了选择。
模型本身决定了上限：从 GPT-6.1 Sol 换到 GPT-6.2 Sol，对哪些题能通过的影响比任何框架手段都大。

\para{范围}本文不提出新模型，不给出单个上下文装不下的任务上的结论，也不声称任何排行榜名次。所有运行都由我们
自己在一台机器上完成，大多通过订阅账号登录，成本是按 API 公开标价重新计算的（第~\ref{sec:threats}~节）。

本文结构如下：第~\ref{sec:related}~节回顾相关工作，第~\ref{sec:system}~节介绍系统，第~\ref{sec:method}~节介绍评估方法；
第~\ref{sec:context}--\ref{sec:tb41}~节报告结果；第~\ref{sec:discussion}~节分析多数手段失效的原因；
第~\ref{sec:threats}~节讨论有效性威胁；第~\ref{sec:conclusion}~节总结全文。

\section{相关工作}
\label{sec:related}

\para{编程智能体的基准与排行榜}SWE-bench 让“解决代码仓库里的 issue”成为检验编程智能体的标准方式
\citep{jimenez2024swebench}。Terminal-Bench 在命令行环境中用困难任务评测智能体，每道题包含说明、容器镜像、测试和
参考解 \citep{merrill2026terminalbench}；它通过 Harbor 运行，Harbor 负责准备任务容器、安装智能体，并以统一格式记录
轨迹 \citep{harbor2026}。SWE Atlas 增加了关于真实代码库的问答，由专家编写的评分细则打分
\citep{raghavendra2026sweatlas}。排行榜现在按“智能体 + 模型”排名：Terminal-Bench 榜单列出维护者自己跑的试验
\citep{tbench2026board}；Artificial Analysis 指数用同一套协议跑多个基准，并在 Terminal-Bench~4.1 上审查每次通过的
尝试，检查是否查找答案、是否修改测试 \citep{aa2026coding}；HAL 在不同基准之间统一框架，并在准确率旁边报告成本
\citep{kapoor2026hal}。\citet{kapoor2024agents} 主张智能体评估应当控制成本、保留留出任务；这两条我们都遵守。

\para{智能体框架}ReAct 让推理与工具调用交替进行 \citep{yao2023react}；SWE-agent 表明智能体与计算机之间的接口和
提示词同样重要 \citep{yang2024sweagent}；Agentless 表明，在解决 issue 这类任务上，一条固定流程也能达到智能体循环的
水平 \citep{xia2024agentless}。服务商自家的智能体，Codex CLI \citep{openai2025codex} 和 Claude Code
\citep{anthropic2025claudecode}，都是可以选用助手的单智能体循环。实践者的建议是选择能解决问题的最简单结构
\citep{schluntz2024effective}。我们最终的配置就是这样一个循环；本文补充的是，在固定模型上，把它和我们搭建的各种团队
结构做对照比较。

\para{多智能体系统}AutoGen、MetaGPT 和 ChatDev 把智能体组织成对话或软件公司式的角色
\citep{wu2024autogen,hong2024metagpt,qian2024chatdev}；多智能体辩论在短答案任务上提高了事实性和推理能力
\citep{du2024debate}；采样加投票的效果随智能体数量增长 \citep{li2024moreagents}。\citet{cemri2025mast} 根据执行轨迹
建立了多智能体失败的分类，分为规格问题、智能体之间不一致、任务验证三类。实践者对团队何时值得用看法不一：
\citet{yan2025dontbuild} 认为上下文分散会让多智能体的决策变得脆弱；\citet{anthropic2025multiagent} 则报告，在材料超出
单个上下文的广泛调研任务上，由领队带并行子智能体更好。我们的结果介于两者之间：在单个上下文装得下的任务上，
同一模型组成的团队增加了成本，却没有提高准确率。

\para{验证、自我纠正与选择}自洽性方法对采样出的多条推理路径投票 \citep{wang2023selfconsistency}；Reflexion 和
Self-Refine 把批评意见反馈给下一次尝试 \citep{shinn2023reflexion,madaan2023selfrefine}；CodeT 用生成的测试挑选程序
\citep{chen2023codet}；当有验证器能挑出正确样本时，重复采样能提高覆盖率 \citep{brown2024monkeys,snell2025scaling}。
\citet{huang2024cannot} 发现，没有外部反馈时，模型无法可靠地纠正自己的推理；LLM 评审还存在位置偏差和偏爱自己输出的
偏差 \citep{zheng2023judge}。我们的阅读者、陪审员、选择者和检查者，都是这些思路在智能体框架中的实例；它们失效的原因
正是 \citeauthor{huang2024cannot} 所说的：模型的第二个副本没有第一个副本缺少的任何反馈。

\para{上下文与缓存}模型对长上下文的利用并不均匀 \citep{liu2024lost}；MemGPT 在有限的上下文中换入换出记忆
\citep{packer2023memgpt}。服务商现在会缓存提示词前缀，并对命中缓存的 token 打折计费 \citep{openai2024caching}，
因此提示词前缀是否稳定就成了一个成本变量。我们的上下文纪元正是围绕这个变量设计的。

\para{评估方法}\citet{miller2024errorbars} 主张模型评估要给出误差范围、做配对比较。用同一个留出集做许多决策会让它
过拟合 \citep{dwork2015reusable}；在一个测试集上调出来的结果，换到新测试集上可能下降 \citep{recht2019imagenet}；
在看到数据之前预先登记分析方案，可以限制分析路径上的随意选择 \citep{nosek2018prereg}。编程智能体还会钻评分器的空子
\citep{metr2025rewardhacking}，这正是 AA 审查每次通过的尝试、我们也扫描自己的尝试的原因。

\section{系统}
\label{sec:system}

The One 在个人电脑或 Linux 服务器上运行语言模型智能体团队。桌面应用和无界面服务端 \code{the-one-server} 共用
同一个后端；本文所有实验运行的都是服务端 \citep{zeroth2026server}。

\para{带生命循环的智能体}一个智能体是 SQLite 里的一条持久记录，而不是一个进程。它的状态先提交，再有任何动作；
被唤醒的智能体从这条记录继续，所以崩溃或重启不会丢失已经决定的工作。一轮只能通过生命周期工具结束（交付、交回、
等待）；模型的回复如果没有调用这类工具，循环就继续；出错时会重试，而不是结束运行。

\para{团队即图}一个团队是一份 YAML 定义：若干带有角色、提示词、工具和模型绑定的智能体，连成一张“谁可以把工作
交给谁”的图。领队可以开一张子图（自己的克隆或一个助手团队），把工作交给另一个节点，或者询问使用者。我们最终保留的
配置是\textbf{团队 7.0}：单个领队直接干活，只在被要求时才委派；更早上线的团队（6.0 到 6.6）加了检查节点和规则，
见第~\ref{sec:structure}~节。

\begin{figure}[t]
\centering
\begin{tikzpicture}[
  font=\small,
  node distance=7mm and 14mm,
  box/.style={draw=ink!70, rounded corners=2pt, align=flush center, inner sep=3pt,
              minimum height=11.5mm, text width=24mm, fill=white},
  ext/.style={box, fill=boxfill},
  arr/.style={-{Stealth[length=2.1mm]}, semithick, draw=ink},
  darr/.style={{Stealth[length=2.1mm]}-{Stealth[length=2.1mm]}, semithick, draw=ink},
  lab/.style={font=\scriptsize, text=ink!85, align=flush center, inner sep=1.5pt}
]
  \node[ext] (harbor) {\textbf{Harbor}\\{\scriptsize 每次尝试}\\{\scriptsize 一个容器}};
  \node[box, right=of harbor] (install) {安装已发布的\\服务端};
  \node[box, right=of install] (run) {\code{the-one-server}\\\code{run}\\{\scriptsize 从 stdin 读取说明}};
  \node[box, below=of run] (lead) {\textbf{领队智能体}\\{\scriptsize 团队 7.0}};
  \node[box] (folder) at (install |- lead) {任务目录\\{\scriptsize shell、文件}};
  \node[ext, right=of lead] (model) {\textbf{模型 API}\\{\scriptsize GPT-6.2 Sol，max}};
  \node[box, below=of lead] (out) {\code{run.json}、\\ATIF 轨迹};
  \node[box] (verify) at (install |- out) {验证器\\{\scriptsize 隐藏测试}};
  \node[ext] (reward) at (harbor |- out) {\textbf{奖励}\\{\scriptsize 0 或 1}};

  \draw[arr] (harbor) -- (install);
  \draw[arr] (install) -- (run);
  \draw[arr] (run) -- (lead);
  \draw[darr] (lead) -- node[lab, above] {模型\\调用} (model);
  \draw[darr] (lead) -- node[lab, above] {工具\\调用} (folder);
  \draw[arr] (lead) -- node[lab, right, align=left, xshift=1.5pt] {交付或\\超时} (out);
  \draw[arr] (out) -- node[lab, above] {然后} (verify);
  \draw[arr] (verify) -- (reward);

  \begin{scope}[on background layer]
    \node[draw=ruleGray, dashed, rounded corners=4pt, inner sep=3.2mm,
          fit=(install) (run) (lead) (folder) (out) (verify)] (ctr) {};
  \end{scope}
  \node[lab, anchor=south west] at (ctr.north west) {任务容器};
\end{tikzpicture}
\caption{通过 Harbor 适配器进行的一次 Terminal-Bench 尝试，即最终评估时的运行方式。在团队 7.0 中，领队自己干活，
只在被要求时才开助手；其余各框对任何通过 Harbor 运行的智能体都一样。智能体自己的记录（\code{run.json}）包含结果、
token 用量和按公开标价计算的成本；奖励只来自验证器。}
\label{fig:flow}
\end{figure}

\para{上下文纪元}一次模型请求由三部分组成：稳定底座（指令、技能、项目笔记）、只追加的日志和一段简短的尾部
（图~\ref{fig:epochs}）。新的上下文在出现的位置追加，之后不再移动，因此服务商的提示词缓存可以持续命中。底座只在
指定的边界重新渲染：更换模型或工具、生成新的摘要，或者漂移超过 40k token。更早的“旧编译器”（classic）每一轮都会
重新渲染指令。

\begin{figure}[t]
\centering
\begin{tikzpicture}[font=\small, x=1mm, y=1mm]
  \node[anchor=east, align=right] at (-2,14) {旧编译器\\[-2pt]{\scriptsize 每一轮}};
  \fill[oursLight!70] (0,10) rectangle (38,18);
  \node at (19,14) {指令，每轮重新渲染};
  \fill[gray!16] (38,10) rectangle (96,18);
  \node at (67,14) {历史};
  \fill[gray!34] (96,10) rectangle (112,18);
  \node at (104,14) {尾部};
  \draw[line width=1.8pt, ours] (0,7.5) -- (5,7.5);
  \node[anchor=west, font=\scriptsize, text=ours!85!black] at (6.5,7.4)
    {缓存常在系统部分就断开（30--40\% 的请求）};
  \node[anchor=east, align=right] at (-2,-6) {纪元};
  \fill[gray!30] (0,-10) rectangle (38,-2);
  \node at (19,-6) {稳定底座};
  \fill[gray!16] (38,-10) rectangle (96,-2);
  \node at (67,-6) {只追加日志};
  \fill[gray!34] (96,-10) rectangle (112,-2);
  \node at (104,-6) {尾部};
  \draw[line width=1.8pt, ink!75] (0,-12.5) -- (96,-12.5);
  \node[anchor=north, font=\scriptsize] at (48,-13.5)
    {缓存前缀：89--94\% 的提示词 token 读自缓存（配对运行，20 道题）};
\end{tikzpicture}
\caption{旧编译器与上下文纪元下的请求布局。纪元在指定边界之间（更换模型或工具、生成新摘要、漂移超过 40k token）
保持底座不变，只做追加，因此尾部之前的内容都能命中缓存。}
\label{fig:epochs}
\end{figure}

\para{摘要}当历史变长时，一个摘要器（“冥想”，界面上显示为“在压缩对话历史”）会用一份精简摘要替换它。当前任务、
对任务的修改和最新的汇报会原样保留，因为不这样做的话，摘要只能保住大约三分之一的规则句（第~\ref{sec:context}~节）。

\para{无界面运行与 Harbor 适配器}\code{the-one-server run} 从标准输入读取任务，在任务目录里运行团队，没有人回答它的
问题，最后写出 \code{run.json}（结果、用量、按公开标价计算的成本、改动的文件）以及每个智能体的 ATIF 轨迹。Harbor 的
安装式智能体在每个任务容器里安装已发布的服务端，然后调用 \code{run}（图~\ref{fig:flow}）。适配器已经公开
\citep{zhang2026adapter}；复现最终运行的命令见附录~\ref{app:repro}。

\section{评估方法}
\label{sec:method}

三个实验项目在同一个产品上先后进行（表~\ref{tab:programs}）。V1 在自建任务集上检验团队设计和效率；V8 通过 Harbor
在 Terminal-Bench~4.0 上检验各种手段；最后一次运行按 Artificial Analysis 跑智能体的方式，把选定的配置放到
Terminal-Bench~4.1 上。

\begin{table}[t]
\centering
\small
\caption{三个实验项目。V8 的试验数包括每一次已评分的尝试，含筛查和重跑。}
\label{tab:programs}
\begin{tabularx}{\textwidth}{@{}l l L >{\raggedright\arraybackslash}p{30mm} >{\raggedright\arraybackslash}p{33mm}@{}}
\toprule
项目 & 日期（2026） & 任务 & 模型 & 规模 \\
\midrule
V1 & 9/28 -- 10/26 & 自建任务集：自编任务、代码仓库任务、公开任务（Terminal-Bench、DeepSWE）和问答任务，都带隐藏
  评分器 & 先 GPT-6 Astra，后 GPT-6.1 Sol，中等强度 & 两个本地评估实例上 120 多批 \\
V8 & 10/27 -- 11/19 & Terminal-Bench~4.0：不需要 GPU 的 63 道题，分为可读 33 道、未读 30 道 & GPT-6.1 Sol、GPT-6.1
  Terra（筛查）、GPT-6.2 Sol & 一台 Linux 电脑上 315 个 Harbor 作业、6,223 次已评分试验 \\
最终 & 11/20 & Terminal-Bench~4.1，AA 的数据集摘要版本，63 道题 & GPT-6.2 Sol，max 强度 & 189 次尝试，每题 3 次 \\
\bottomrule
\end{tabularx}
\end{table}

\subsection{任务集}

\para{V1 任务集}每道题都带一个隐藏评分器和一份参考解。一道题进入任务集之前，先由一个全新的模型在空白工作区里盲做
一遍，用它的得分检验评分器；在早期的一批里，这样找出了两个真实的评分器缺陷。已经饱和或不公平的题会被撤下
（10 月 9 日撤下 52 道，剩 32 道在用）。一个设计在它所针对的题上胜出之后，还必须在不针对的题上站得住，才能上线。

\para{V8 的划分}10 月 30 日，63 道能运行的 Terminal-Bench~4.0 题目被分成 25 道开发题和 38 道留出题。第一次读取留出
结果（11 月 4 日）之后，榜单上有智能体胜过我们的 8 道留出题被打开研究。这样剩下 33 道\textbf{可读}题，我们阅读它们的
轨迹、失败的测试和题目文件来设计手段；另有 30 道\textbf{未读}题用于确认。未读题只以总和的形式出现：奖励、成本、
分钟数和模型调用次数的合计；它们的题名、轨迹和测试从未打开过，当时没有，写这篇论文时也没有。

\subsection{基础设施}

V8 使用 Harbor 0.24.0（最终运行用 0.26.0）\citep{harbor2026}，运行在一台 32 线程、96\,GB 内存的台式电脑上；在更高
负载下一次内存交换风暴损失了 44 次运行，此后同时运行约 30 次试验。团队自身的时间上限是 465 分钟，题目的上限是
480 分钟，因为服务端默认的 3 小时曾把一道需要 8 小时的题提前截断。11 月 5 日之前，运行在一台 Mac 上的 Docker
虚拟机里进行，每次同时跑两次试验；它的留出基线比那台电脑低，这也是每次比较都在同一批次、同一台机器上进行的原因
之一。

\subsection{指标}

题目 $t$ 的第 $d$ 次运行的奖励 $r_{t,d}$ 为 0 或 1。我们的得分是各题平均奖励之和：
\begin{equation}
  S = \sum_{t \in T} \frac{1}{n_t} \sum_{d=1}^{n_t} r_{t,d},
  \label{eq:score}
\end{equation}
因此 $S/|T|$ 就是通过率。缺失或没有奖励的尝试记为 0。成本是该次运行的 token 按服务端价格表中 OpenAI 的 API 公开标价
（2026 年 10 月 18 日核对）计算的，尽管多数运行实际使用的是 ChatGPT（Codex）账号登录。时间是智能体自己每次尝试所用的
时间；V1 的效率研究用的是到最后一次模型调用为止的活跃时间。

\subsection{统计分析}

在 Terminal-Bench 题集上，实验组和基线在每道题上各运行两次或更多。我们检验 $S$ 的差异时，在每道题内部把各次运行
在实验组和基线之间重新分配（20,000 次置换），由此得到零假设下的标准差和 $p$ 值。在未读的 30 道题上每边两次运行时，
零假设标准差约为 1.4（各次读取在 1.33 到 1.53 之间），因此一次读取能分辨约 2.7 道题的差异；每边四次运行时约为 0.94。
凡给出标准误的地方，标准误都是按题计算的。V1 对配对运行使用成本或时间对数比的精确 Wilcoxon 符号秩检验
\citep{wilcoxon1945}。

\subsection{实验规范}

\begin{itemize}[leftmargin=1.4em]
  \item \textbf{同批基线。}每个实验组都和它的基线在同一时段运行，因为曾有一个基线在两批之间移动了 +1.67。
  \item \textbf{空闲重评。}每个零分都在空闲机器上重评一次，因为负载下验证器超时会产生假零分。在空闲机器上这种情况
    很少：两次大规模重评中，43 个零分里 43 个、54 个零分里 54 个仍为零；第三次重评中，74 个里有 1 个变了。
  \item \textbf{预先登记。}每批实验之前，都把规则、分支和我们自己的预测写进实验记录本 \citep{nosek2018prereg}；V8 的
    记录本最终有 9,276 行、387 节。我们还请一个独立的审查智能体专门找设计和解读中的漏洞。
\end{itemize}

\section{上下文与成本（RQ1）}
\label{sec:context}

让每个智能体的提示词只追加，使每次运行的成本降低了三分之一到一半，得分不变。之后测量的上下文手段让成本变化了
7\% 到 20\%，没有一个改变准确率（表~\ref{tab:levers}）。

\para{只追加的提示词}旧编译器每一轮都重新渲染指令，因此 30\% 到 40\% 的请求在系统部分就断开了缓存。在 9 月 30 日的
一次配对运行中（同一次重启、同样负载，20 次运行，全部解出），纪元方式有 89\% 到 94\% 的提示词 token 读自缓存，旧方式
为 76\% 到 81\%。旧方式每次运行的成本高出 50\% 到 108\%，例如一道单人 SQL 题是 0.81 美元对 0.44 美元。纪元边界之外的
前缀改写只占 9,897 次请求的 0.25\%；空闲长达一小时之后，缓存命中率仍有 98\% 到 99\%。

\para{服务商的缓存需要正确的请求头}在 ChatGPT 后端上，只带缓存键什么也缓存不了。加上会话和线程请求头之后，命中率
升到 97.8\% 到 99.6\%；重放轮次状态又让一轮中途的请求从 0 升到 98.5\%。在这些修复之前，一个长期运行的团队一周花掉了
5.55 亿个提示词 token：它在 872k token 窗口的 85\% 处才压缩，而后端并不支持这么大的窗口；它的提示词涨到 744k token，
而且完全没有命中缓存。现在压缩在 231k token 时触发。

\para{委派的代价在冷启动}在第 2 批（9 月 29 日）中，委派的运行快了 25\%，但发送的提示词 token 是原来的 2.46 倍，
未命中缓存的 token 是 7.57 倍。52 个被派生的智能体，每个的第一次请求都约有 42k token 未命中缓存，占这些运行未命中
token 的 57\%。随后做了两项修复：智能体自己的笔记改为放在共享指令之后、作为用户消息发送（一个审查者第一次请求中
命中缓存的 token 从 42k 里的 18k 升到 38k）；克隆体继承源智能体的工具和纪元（审查者第一次请求的缓存命中率从 13\% 升到
96\%）。再加上团队共用的缓存键、静态的交接工具和一个预热关口，每次委派运行的未命中 token 从 643k 降到 390k。

\para{延迟跟着 token 走，不论是否命中缓存}在同一次运行、同一时刻发出的 2,354 对匹配请求上，首个事件的到达时间每多
10k 个命中缓存的 token 增长到 1.09 倍，每多 10k 个未命中的 token 增长到 1.11 倍；消息条数没有影响。领队是关键路径：
一次运行约 90\% 的活跃时间花在它自己的模型调用上，它 86\% 的轮次只调用一个工具。减少领队的轮次能省时间；缩小
提示词主要省钱。

\begin{table}[t]
\centering
\footnotesize
\caption{V1 任务集上的上下文手段。比值为实验组除以同批基线；“完全解出”指拿到满分的运行。}
\label{tab:levers}
\begin{tabularx}{\textwidth}{@{}>{\raggedright\arraybackslash}p{36mm} >{\raggedright\arraybackslash}p{18mm} L L l@{}}
\toprule
手段 & 批次与规模 & 成本或 token（对基线） & 得分与时间 & 结论 \\
\midrule
只追加的纪元（对旧编译器） & 9/30，20 次运行 & 成本低 33\% 到 52\% & 两边全部解出 & 默认开启 \\
提前写项目笔记（播种运行） & b101，5 道题 & token 0.84 & 时间 0.83；完全解出 5 对 4 & 采用 \\
读取指引（不整文件加载） & b96，24 对 & 未命中 0.80（$p = 0.001$），成本 0.86 & 时间 0.95；完全解出 18 对 16 &
  采用 \\
从未调用的工具只给简短 schema & b97，24 对 & 领队首次请求 28.9k 降到 21.9k token；token 0.89，成本 0.93 & 时间 0.97；
  完全解出 18 对 15 & 采用 \\
精简工具结果 & b100，24 对 & token 0.98，成本 0.98 & 时间 1.06 & 放弃 \\
多个任务共用一个会话 & s1，5 道题 & 成本 1.06 & 得分相同 & 放弃 \\
60k token 时移出旧工具结果 & b95，24 对 & 未命中 1.29，成本 1.14 & 时间 1.02 & 放弃 \\
\bottomrule
\end{tabularx}
\end{table}

\para{摘要会丢规则}对 861 份已提交摘要的审计发现，摘要长度中位数只有原文的 1.1\%，规则句的 token 只有 36.9\% 被原样保留；
每一份后续摘要大约只保住前一份所保留内容的一半。现在摘要器会原样保留当前任务、对任务的修改和最新的汇报。对话回忆
测试显示不出这项改动的收益：每个实验组在三次摘要之后都记住了全部八条植入的事实，所以这项改动的依据只有那次审计。

\para{RQ1 的回答}上下文管理是一个成本手段。让提示词前缀稳定，比之后所有削减 token 的做法加起来都值；而一个打断
缓存的削减，花掉的比省下的多（第~\ref{sec:discussion}~节）。至于时间，要紧的是领队的轮次，而不是提示词的大小。

\section{团队结构与运行时干预（RQ2）}
\label{sec:structure}

在单个上下文装得下的任务上，一个能力强的负责者不输于我们搭建的任何团队，成本只是团队的一小部分。增加的智能体在
它们所针对的题上胜出，在不针对的题上持平或落后，成本通常是两到五倍。

\para{小任务不委派，大任务才委派}在第一批任务集上（9 月 28 日），默认团队从不委派，单个智能体同样好，而且略便宜
（10.9 分钟对 12.1 分钟，1.24M token 对 1.41M）。在更大的多模块任务上（9 月 29 日），领队在其中两道题的全部六次运行里
都做了委派，在调研类任务上一次也没有；36 次运行全部解出，所以结构只能影响时间和成本。

\para{为难题搭建的结构不能迁移}从 10 月 8 日起，我们针对上线团队做不出的难题，每天搭一种结构（表~\ref{tab:structures}）。
对这些失败的诊断发现，每一次都是对需求中一句含糊或没说清的话的误读，而且多数运行都以同样方式误读，因此投票和
增加审查者几乎无从下手。在每个节点上加阅读规则，到了样本外反而更差，因为这些规则让模型觉得可以跳过文字没有
明说的部分。

\begin{table}[t]
\centering
\footnotesize
\caption{V1 中为难题搭建的团队结构，每个都与同批基线比较（实验组对基线，通过的运行数）。设计题是该结构所针对的题；
留出题没有用于搭建它。}
\label{tab:structures}
\begin{tabularx}{\textwidth}{@{}>{\raggedright\arraybackslash}p{34mm} L L >{\raggedright\arraybackslash}p{21mm} >{\raggedright\arraybackslash}p{27mm}@{}}
\toprule
结构 & 设计题 & 留出题 & 每次运行成本 & 结论 \\
\midrule
每个节点加阅读规则 & 9/15 对 4/15 & 2/10 对 4/10 & 约 1.00 美元对 0.76 美元 & 放弃：规则让工作缩水 \\
三个构建者加一个证据评审（集成） & 15/24 对 4/15；某道题第一次被解出 & 3/10 对 8/20 & 2.85 美元对 0.53 美元 & 放弃 \\
三个阅读者加一个调和者 & 4/10 对 4/21 & 1/9 对 11/28 & 4.01 美元对 1.31 美元 & 放弃 \\
与构建同时启动的验收检查者 & 3/5 对 2/15（细节密集的题） & 8/10 对 7/10（新题集） & 2.03 美元对 1.58 美元 &
  作为关口采用，只用于规格类需求 \\
画类型化连线图的架构师 & 0/6；五次完全打平 & 未运行 & 未记录 & 放弃 \\
\bottomrule
\end{tabularx}
\end{table}

\para{V1 上线了什么}在留出题上与基线持平的单个检查者确实上线了，一次加一个关口：验收（团队 6.0）、针对缺陷报告的
症状检查（6.1）、代码仓库惯例（6.2）、检查者交回时的定夺步骤（6.3；在其目标题上 17/48 对 10/48）以及组合测试者
（6.4）。更重的变体全都没通过留出检验，包括双检查者版本（3/9 对 5/9）和定夺者节点（合计 20/40 对 21/40）。10 月 18 日对
120 次运行中 61 次失败的普查发现：33 次是误读需求，16 次是自测比隐藏测试宽松，12 次是领域知识缺口。

\para{单领队在 Terminal-Bench 上同样胜出}在 V8 中，团队 7.0（单个领队直接干活，只在被要求时委派）以 max 强度运行，
在未读的 30 道题上比之前中等强度的领队高 $+3.17$（标准误 1.39），10 道题上升、3 道下降，没有丢掉任何一道原本总能
通过的题，成本是两倍。这是两个配置之间的比较：领队的人设和推理强度一起变了，而这一对组合正是我们上线的
（单独比较推理强度见第~\ref{sec:judgment}~节）。按 Codex 自己的协作协议做主动委派（领队不等要求就派生助手智能体），
只在 24 次目标试验中的 2 次上成功，丢了两道守卫题，成本是二到四倍；它的被动形式，即只在被要求时委派，就是团队 7.0
采用的方式。一个检查型团队与单领队持平，成本是两倍。

\para{同模型的第二意见在样本外无效}在动手之前先研究题目的阅读者，在未读题上合计为 $-0.8$（$\pm 0.7$）。在两份已完成
的构建之间做选择的选择者，无论用 GPT 还是 Claude，正确率都在 55\% 左右：失败的那份有明显缺陷时几乎总能选对
（Claude Opus 23 次中 22 次），在需要判断的点上则很少选对（28 次中 5 次）。一个同模型陪审团在 36 对“通过/失败”构建上
判对 12 次、判错 11 次、弃权 13 次。由领队启动的独立评分者在目标试验上通过 3/10，基线为 5/10；它给出的 10 个
“全部通过”结论中有 6 个落在失败的试验上。一个带关口的推理检查者以 12/20 打平。

\para{框架只在回应“通过的运行从不出现的状态”时有用}无界面运行中没有人回答交回。当一次运行很早就交回失败时，框架
回答“你自己决定，继续做”（GO）；有一道题从 14 次中 0 次通过变成 6 次中 6 次，GO 因此上线。在每次首次交付时触发的
检查没有上线。GPT-6.1 上的 HD 在未读题上得 $+1.54$（$\pm 1.31$），成本 2.09 倍，低于预先登记的 $+2.0$ 门槛。O2 是同一
思路针对 GPT-6.2 的措辞，在多次运行结果不一致的可读题上通过 22/28 对 17/28，它的第二轮确实找出了真缺陷；但在未读题上
少了 2.0（15.5 对 17.5，$p = 0.28$），并在 60 次试验中的 50 次里返工了已交付的工作。全局触发的手段在一个失败类型里
得到的收益，要用原本已经做对的题来偿还。

\para{RQ2 的回答}在单上下文任务上，同一模型组成的团队增加了成本，没有提高准确率。有用的东西都很窄：一个只用于
规格类需求的关口，以及在运行时回应一种通过的运行从不会进入的状态（很早就交回失败）。

\section{模型及其判断（RQ3）}
\label{sec:judgment}

模型决定了上限。从 GPT-6.1 Sol 换到 GPT-6.2 Sol，对哪些题能通过的影响比任何框架手段都大；我们写的任何文字都没有改变
模型在决策点上的选择，只有摆在它面前的事实做到了。

\para{多数失败是自信的交付}在可读题上，基线的 227 次失败中有 213 次是模型认为已经完成的交付。每一次都建立在模型
最初几分钟就定下的一种需求解读上，之后它用自己的解读来检验自己，所以它自己的检查都能通过。把同样的题与 Codex
公开的轨迹并排对比：Codex 经过一次长时间思考后定下更强的设计，再去加固它；而我们的运行在自己的检查通过后就停了。
可读题上与 Codex 的全部差距，集中在 Codex 稳定通过的 8 道题上（期望通过数 6.75 对 2.1），没有共同的机制。

\subsection{什么改变了选择}

我们截取真实试验中某个决策点的请求，改动其中一段文字后多次重放（关键节点重放）。重放与原始提示词完全一致，98\%
命中缓存，每个样本约需 13 秒。有三类输入改变了选择：
\begin{itemize}[leftmargin=1.4em]
  \item \textbf{在恰当时刻、以用户轮次给出的事实或指令。}在题目里加一行，指明某个缺陷已经确定，使下一步在 24 次重放
    中有 23 次写出了无偏估计量。第~\ref{sec:structure}~节的 GO 回答就属于这一类。
  \item \textbf{拿掉误导性的线索。}在一个人为设置的选择探针中，告诉模型代码文档里写的是哪个选项时，它在 16 次重放中
    16 次都保留了那个选项；不告诉它时，16 次都选对了。
  \item \textbf{把数据的测量结果摆到它面前。}在一道题上，这让两个 GPT 模型都从 8 次中 0 次变成 8 次中 8 次。
\end{itemize}

\subsection{什么没有改变选择}

\begin{itemize}[leftmargin=1.4em]
  \item \textbf{指令或人设中的建议}，试了五种措辞。“注释不是证据”这句话在 24 次决策中一次也没有改变；为 GPT-6.2 写的
    一段人设也没有改变任何可测的行为。
  \item \textbf{同一模型的预备步骤}，在开工前列出并决定待定的点。只对一道题起作用。
  \item \textbf{在领队旁边放另一个模型做通用阅读者或选择者。}Claude Opus 做阅读者 8 次中 0 次有效，而且起了反作用；
    在两份构建之间做选择，正确率为 0.55。另一个模型只在它知道 GPT 不知道的东西时才有帮助：在 V1 中，一个 Opus 定夺者
    把一道统计题修好了 4 次中 4 次，基线为 4 次中 0 次。
  \item \textbf{实时送达的提示。}一个适配器钩子把一段通用提示放在运行中第一次写文件的结果之后，并在之后每次请求的
    同一位置出现，因此保住了提示词缓存。作为用户轮次时，这段提示让下一步在两个截取的决策点上分别有 8 次中 8 次、8 次中
    5 到 6 次去检验关键点；作为工具结果的后缀时，分别是 8 次中 5 到 6 次和 8 次中 0 次；不加提示时 4 次中 0 次。然而在
    完整运行中，可读的 33 道题得分 21.0 对 21.0，成本相同：在它针对的 7 道开放解读型题上 $+1.0$，在其余 26 道上 $-1.0$。
    在一道题上，模型做了被要求的检查，打印出评分器期望的那种形式的分数（3.239，它自己的形式是 3.240），然后仍保留了
    自己的形式。
\end{itemize}

在截取的决策点上重放，会高估完整运行的效果：一段“先探查”的文字让一个决策从 24 次中 0 次变成 24 次中 24 次，但对完整
运行的影响是 $-0.04$（$\pm 0.85$）。恰当时刻的提示改变的是过程，不是判断；判断跟随模型自己的先验。我们因此不再把
建议当作在这些题上影响 GPT 判断的手段。

\subsection{模型版本与推理强度}

\begin{table}[t]
\centering
\small
\caption{团队 7.0、max 强度下，GPT-6.1 Sol 与 GPT-6.2 Sol 在 Terminal-Bench~4.0 上的比较。$S$ 为各题平均奖励之和
（式~\eqref{eq:score}）；“$N$ 次取最好”假设有一个完美的选择者从 $N$ 次运行中挑选。}
\label{tab:models}
\begin{tabular}{@{}lll@{}}
\toprule
指标 & GPT-6.1 Sol & GPT-6.2 Sol \\
\midrule
可读 33 道，$S$ & 约 18.8 & 约 21 \\
未读 30 道，同一批次，$S$ & 20.5 & 18.0（差 $-2.5$，$p = 0.15$） \\
未读 30 道，单次运行的平均 & 19.25（8 次运行） & 17.75（4 次运行） \\
未读 30 道，2 次取最好 / 4 次取最好 & 21.04 / 21.77 & 19.33 / 20.0 \\
多次运行结果不一致的未读题 & 8 次运行中 9/30 & 4 次运行中 6/30 \\
每次试验的模型调用次数 & 46.9 & 31.4 \\
每次试验的智能体用时（分钟，未读） & 33.7 & 28.8 \\
每次试验按公开标价的成本（未读） & 1.16 美元 & 1.15 美元 \\
\bottomrule
\end{tabular}
\end{table}

\para{从 GPT-6.1 Sol 到 GPT-6.2 Sol}表~\ref{tab:models} 在同一配置下比较了两个模型。GPT-6.2 的模型调用少三分之一，
用时少 15\%，按公开标价的单次成本相同。它能通过 GPT-6.1 很少通过的开放解读型题（某道题 5 次中 5 次，GPT-6.1 为 11 次中
0 次），却从不通过一些 GPT-6.1 有时能通过的题（8 次中 0 次，GPT-6.1 为 50 次中 17 次和 50 次中 16 次）。它也更确定：它的
瓶颈是每次运行都以同样方式做出的单个方法选择，所以多跑几次对它的帮助更小。我们在最终配置中保留了 GPT-6.2，因为未读题
上的差异在噪声范围内，而我们比较的对象用的也是 GPT-6.2。

\para{推理强度}在 V1 中，中等强度在十道缺陷题上与高强度相当（28/30 对 29/30），成本低 38\%，但在多缺陷题上落后
（0/6 对 3/6，$p \approx 0.06$）。在一组可读的 Terminal-Bench 筛查题上，之前的领队在 max 强度下 38 次试验通过 19 次，
中等强度 21 次通过 6 次，高强度 15 次通过 3 次，所以“高”并不是便宜的折中。对单领队而言，前 12 次调用用 max、之后用
中等强度，通过 15 次中 5 次，全程 max 为 15 次中 11 次，成本相同。max 强度因此成为编程任务上线时的深度，第~\ref{sec:structure}~节
的配置比较就是它的留出检验。

\para{RQ3 的回答}框架能改变模型下一步做什么（做一次检查、继续做），但改变不了它在选择点上的结论；只有摆到它面前的
事实和测量结果改变了选择。剩下的失败都是自信的单一解读，所以能动到它们的手段必须提供证据，而不是建议。

\section{与人协作}
\label{sec:person}

基准测试只把一个完整的任务交给智能体一次；而一个真实的人会纠正它、过后再回来，并期望它记得。四个产品测试台测量了
这一面（表~\ref{tab:benches}）。它们的结果与编程实验指向同一个方向：能传下去的，是写在下一个智能体会读到的地方的
东西，而不是存在记忆工具里的东西。

\begin{table}[t]
\centering
\footnotesize
\caption{产品与人协作的测试台。由脚本扮演的人和小组代替真人；“之前”和“之后”指最后一列所列改动之前和之后的产品。}
\label{tab:benches}
\begin{tabularx}{\textwidth}{@{}>{\raggedright\arraybackslash}p{20mm} >{\raggedright\arraybackslash}p{38mm} L >{\raggedright\arraybackslash}p{24mm}@{}}
\toprule
测试台 & 设置 & 结果 & 上线的改动 \\
\midrule
验收习惯（11/17--18） & 一个脚本扮演的人，有四条任何任务都没写明的代码习惯；每个序列六道 Python 题；过程中给出纠正 &
  之前：一条被纠正过的习惯，经由主助手时 51 次中 0 次再被违反（它把这个人的原话贴进每份任务简报），直接交给团队时
  51 次中 51 次被违反。在留出的人上：团队违反被纠正习惯 14/96，主助手 0/48；一次性要求被误当成规则保留 0/108 和
  3/81 & 常设规则文件和一个保存规则的工具 \\
记忆教训（11/18） & 29 个序列，后面的任务需要用到前面任务的教训 & 全新运行 0/6，告知后 4/4，团队 6/8，带记忆的团队 5/8，
  主助手 8/8；团队恰好在读了前一任务目录时通过（11/11 对 0/5）；压缩之后 12/12 & 不重写记忆；教训通过文件和任务简报
  传递 \\
跟进（10/25） & 脚本化的小组把 7 种回测变体之一交给主助手；每组 21 次运行 & 跟进率前后都是 18/21；评分细则得分从 0.49
  升到 0.88 & 主助手的交回文字 \\
决策模型（10/25--26） & 盲审阅读者判断一个小组是否揭示了脚本设定的决策模型 & 有实验设计技能时 24/24 对 16/24
  （$p = 0.004$），token 为 0.88 倍 & 实验设计技能 \\
\bottomrule
\end{tabularx}
\end{table}

有两个结果改变了设计。主助手之所以守住每条习惯，是因为它把这个人的原话改写进了每份任务简报；团队只有在这些话进入
它们会读的文件时才守得住。并且在 104 个序列任务中，没有一个智能体主动调用过记忆工具；所以常设规则现在放在下一个
智能体会读的文件里，教训则留给工作本身的文件和主助手的任务简报去传递。

有一条工作线没有找到值得上线的东西。“进化实验室”跑了八轮（模型成本 59.55 美元），追查一次失败，结果发现那次失败
主要来自我们自己写的任务摘要：摘要改正之后，上线团队 4 次中 4 次解出该题，之前是 21 次中 3 次。

\section{最终配置在 Terminal-Bench 4.1 上的表现（RQ4）}
\label{sec:tb41}

The One 0.5.2 的团队 7.0，用 GPT-6.2 Sol、max 强度，在能运行的 63 道题上得分 60.9\%，在全部 66 道题上（3 道 GPU 题记为
零分）得分 58.1\%。这与 AA 跑出的同一模型、同一强度下的 Codex（58.1\%）持平，与 Terminal-Bench 维护者跑出的结果（61.5\%）
也在噪声范围之内。

\para{设置}Harbor 0.26.0，使用 AA 的数据集摘要版本；每个任务容器里安装已发布的服务端 0.5.2；GPT-6.2 Sol、max 强度，
通过 ChatGPT 账号登录；每题三次尝试，与 AA 的协议一致；在一台没有 GPU 的 Linux 电脑上同时运行 30 次；遵守
Terminal-Bench 8 小时的智能体时限。有两次留出题的尝试在智能体启动前下载服务端时中断；按 Terminal-Bench 维护者重跑
基础设施错误的做法，我们重跑了这两次，结果仍为 0 分。与我们的 Terminal-Bench~4.0 副本相比，4.1 只有一道题不同。

\begin{table}[t]
\centering
\small
\caption{The One 0.5.2、团队 7.0、GPT-6.2 Sol、max 强度在 Terminal-Bench~4.1 上的结果，每题三次尝试。在全部 66 道题上
（3 道 GPU 题记为 0），平均奖励为 58.1\%（115/198）。}
\label{tab:tb41}
\begin{tabular}{@{}lrrr@{}}
\toprule
 & 可读 33 道 & 留出 30 道 & 全部 63 道 \\
\midrule
平均奖励 & 62.6\%（62/99） & 58.9\%（53/90） & 60.9\%（115/189） \\
至少解出一次的题数 & 24 & 20 & 44 \\
每次尝试的智能体用时（分钟），中位数（均值） & 24.7（34.5） & 18.1（20.5） & 20.5（27.8） \\
每次尝试按公开标价的平均成本 & 1.61 美元 & 1.06 美元 & 1.35 美元 \\
\bottomrule
\end{tabular}
\end{table}

表~\ref{tab:tb41} 按划分给出结果；留出题的得分与可读题相差不到四个百分点。整次运行按公开标价花费 255 美元，墙钟时间
6.97 小时。一次尝试平均使用 316 万个输入 token，其中 90\% 命中缓存；输出 4.5 万个 token，其中 55\% 是推理 token。
有一道可读题的三次尝试每次都用了约五个小时，拉高了平均用时。

\begin{figure}[t]
\centering
\begin{tikzpicture}
\begin{axis}[
  width=0.50\textwidth, height=7.6cm,
  xmin=45, xmax=71, ymin=0.4, ymax=11.6, y dir=reverse,
  xtick={45,50,55,60,65,70},
  ytick={1,...,11},
  yticklabels={
    {Kimi Code + Kimi K3.1 Pro},
    {Claude Code + Claude Sonnet 5.5 (max)},
    {Antigravity + Gemini 4 Neon},
    {Claude Code + Claude Fable 5.2 (xhigh)},
    {Codex + GPT-6.2 Astra (max)},
    {Codex + GPT-6.2 Astra (xhigh)},
    {\textcolor{ours}{\textbf{The One + GPT-6.2 Sol (max)，63 道}}},
    {Codex + GPT-6.2 Sol (max)},
    {\textcolor{ours}{\textbf{The One + GPT-6.2 Sol (max)，66 道}}},
    {Codex + GPT-6.2 Sol (xhigh)},
    {Codex + GPT-6.2 Sol (medium)}},
  yticklabel style={font=\footnotesize},
  tick label style={font=\footnotesize},
  xlabel={三次尝试的平均奖励（\%）},
  xlabel style={font=\footnotesize},
  xmajorgrids, grid style={ruleGray!45},
  axis x line*=bottom, axis y line*=left,
  y axis line style={draw=none}, ytick style={draw=none},
  clip=false,
]
  \addplot[xcomb, draw=ruleGray, line width=0.6pt] coordinates
    {(67.2,1) (65.7,2) (65.2,3) (64.1,4) (63.6,5) (62.1,6) (60.9,7) (58.1,8) (58.1,9) (57.6,10) (49.5,11)};
  \addplot[only marks, mark=*, mark size=2.4pt, color=other,
    nodes near coords, nodes near coords style={font=\scriptsize, anchor=west, xshift=1pt, text=ink!85},
    point meta=explicit symbolic] coordinates
    {(67.2,1) [67.2] (65.7,2) [65.7] (65.2,3) [65.2] (64.1,4) [64.1] (63.6,5) [63.6]
     (62.1,6) [62.1] (58.1,8) [58.1] (57.6,10) [57.6] (49.5,11) [49.5]};
  \addplot[only marks, mark=*, mark size=2.6pt, color=ours,
    nodes near coords, nodes near coords style={font=\scriptsize\bfseries, anchor=west, xshift=1pt, text=ours},
    point meta=explicit symbolic] coordinates {(58.1,9) [58.1]};
  \addplot[only marks, mark=o, mark size=2.6pt, color=ours, thick,
    nodes near coords, nodes near coords style={font=\scriptsize\bfseries, anchor=west, xshift=1pt, text=ours},
    point meta=explicit symbolic] coordinates {(60.9,7) [60.9]};
\end{axis}
\end{tikzpicture}
\caption{AA 协议下的 Terminal-Bench~4.1（每题三次尝试）。灰色：AA 自己跑的其他智能体，2026 年 11 月 20 日读自其编程智能体
页面 \citep{aa2026coding}。橙色实心：我们在全部 66 道题上的结果，3 道 GPU 题记为 0，可与 AA 各行直接比较。空心：我们在
能运行的 63 道题上的结果，不可直接比较。横轴从 45\% 开始。}
\label{fig:board}
\end{figure}

\para{与榜单比较}AA 按基准给出 token 数，但成本和时间只按智能体、在它的全部基准上汇总给出，因此我们按同样的公开标价
自己计算了它的 Codex 各行（表~\ref{tab:cost}）。Codex 在 GPT-6.2 Sol、max 强度下每次尝试用 448 万输入 token 和 5.1 万输出
token，按公开标价约 1.42 美元；我们少用 30\% 的输入 token，花 1.35 美元。Codex 在 GPT-6.2 Astra、max 强度下得分 63.6\%，
每次尝试约 5.04 美元。维护者自己的 4.1 榜单在 66 道题上各跑五次，列出 Codex + GPT-6.2 Sol（max）为 61.5\%，Codex +
GPT-6.2 Astra（max）为 62.1\% \citep{tbench2026board}。两个榜单的头部都是 Claude、Gemini 和 Kimi 的智能体，以及一个多模型
智能体：GitHub Copilot 的 Copilot Fusion（在维护者榜单上为 71.2\%）。

\begin{table}[t]
\centering
\small
\caption{每次 Terminal-Bench~4.1 尝试的 token 与公开标价成本。我们这一行是实测值；Codex 各行由 AA 给出的每次尝试
Terminal-Bench token 数和按智能体汇总的缓存命中率（在其全部基准上统计）按 OpenAI 公开标价计算（GPT-6.2 Sol：
每百万未命中输入、命中输入和输出 token 分别为 2、0.10 和 10 美元；GPT-6.2 Astra：8、0.40 和 40 美元）。}
\label{tab:cost}
\begin{tabular}{@{}lrrrrr@{}}
\toprule
智能体与模型（强度） & 66 道题得分 & 输入 & 命中缓存 & 输出 & 成本 \\
\midrule
\textbf{The One} + GPT-6.2 Sol (max) & 58.1\% & 3.16M & 90\% & 45.2k & \$1.35 \\
Codex + GPT-6.2 Sol (max) & 58.1\% & 4.48M & 95\% & 51.4k & $\approx$\$1.42 \\
Codex + GPT-6.2 Astra (max) & 63.6\% & 4.06M & 95\% & 50.0k & $\approx$\$5.04 \\
\bottomrule
\end{tabular}
\end{table}

\para{此前在 Terminal-Bench~4.0 上}在同样的划分下，该配置用 GPT-6.1 Sol 得分约为 63 道中的 38 道，用 GPT-6.2 Sol 约为
39 道；榜单上公开的试验中，Codex 0.159 + GPT-6.1 Sol 为 36.4，Claude Code 2.1.284 + Claude Opus 5.5 为 39.8。在未读的
一半上，我们的 GPT-6.1 配置比 Codex 高 $+1.23$（六次运行的 95\% 区间为 $-0.30$ 到 $+2.77$）。Claude Code 在可读一半上的
领先，来自 7 道开放解读型题，它在这些题上 35 次通过 34 次。按公开标价，我们每次试验在未读题上花 1.37 美元，Codex
1.51 美元，Claude Code 11.39 美元。第一次完整的基线（10 月 31 日，当时上线的团队、中等强度、在 Mac 上）通过了 63 道中的
29 道。

\para{SWE-Atlas-QnA}在 124 道题的问答集上 \citep{raghavendra2026sweatlas}（10 月 29 日，GPT-6.1 Sol、中等强度，当时的团队），
AA 的评分细则评分器判定满足 58\%（已评分的 120 道中 70 道），Scale 的评分细则评分器为 52\%（120 道中 62 道）；每题用时
中位数 9.5 分钟、成本 0.42 美元；当时 AA 列出 Codex + GPT-6.1 Sol（max）为 60\%。在 101 条未满足的评分项中，有 44 条要求
的是运行软件得到的证据。一个问题检查者（$+5$，预先登记的门槛是 $+6$，成本 1.45 倍）和一个高强度领队（$-2$，成本
1.57 倍）都没有保留。

\para{奖励投机扫描}AA 的 4.1 审查会阅读每次通过的尝试，把查找答案、修改测试和访问参考解的尝试记为 0 分
\citep{aa2026coding}。我们的扫描器检查每次尝试的轨迹：是否有提到该基准、某道题或其上游来源的网络调用，以及是否读取、
列出或写入评分路径（测试、参考解、验证器日志、奖励文件）。它在所有尝试中都没有发现这样的网络调用，也没有发现对评分
路径的读取、列出或写入。在留出一侧，某道题的两次通过的尝试写入的文件内容提到了验证器的日志目录；这些轨迹属于留出集，
因此我们没有做判断。

\para{RQ4 的回答}在一个固定的前沿模型上，框架能追上服务商自家的智能体，并在 token 上更省，但还没有超过它。

\section{多数手段为何失效}
\label{sec:discussion}

我们尝试过的大多数东西都没能在留出题上站住（图~\ref{fig:transfer}）。每一种为难题搭建的结构，到了留出题上都失去了
领先；只有验收检查者保持领先，而且它只作为规格类需求的关口上线。这些失败有五个共同原因。我们把它们写出来，是因为
每一个都花了我们几周时间，也会让别人花同样的时间。

\begin{figure}[t]
\centering
\begin{tikzpicture}
\begin{axis}[
  width=0.50\textwidth, height=5.4cm,
  xmin=-35, xmax=55, ymin=0.4, ymax=5.6, y dir=reverse,
  xtick={-20,0,20,40}, minor xtick={-30,-10,10,30,50}, xminorgrids,
  minor grid style={ruleGray!25},
  ytick={1,...,5},
  yticklabels={
    {每个节点加阅读规则（V1）},
    {三个构建者加一个评审（V1）},
    {三个阅读者加一个调和者（V1）},
    {验收检查者（V1）},
    {首次交付时检查，O2（V8）}},
  yticklabel style={font=\footnotesize},
  tick label style={font=\footnotesize},
  xlabel={实验组减同批基线（百分点）},
  xlabel style={font=\footnotesize},
  xmajorgrids, grid style={ruleGray!45},
  axis x line*=bottom, axis y line*=left,
  y axis line style={draw=none}, ytick style={draw=none},
  legend style={font=\footnotesize, draw=none, fill=none, at={(0.5,1.03)}, anchor=south,
    legend columns=2, /tikz/every even column/.append style={column sep=8pt}},
  clip=false,
]
  \draw[ink!55, dashed] (axis cs:0,0.4) -- (axis cs:0,5.6);
  \addplot[draw=ruleGray, line width=1.3pt, forget plot] coordinates {(33.3,1) (-20.0,1)};
  \addplot[draw=ruleGray, line width=1.3pt, forget plot] coordinates {(35.8,2) (-10.0,2)};
  \addplot[draw=ruleGray, line width=1.3pt, forget plot] coordinates {(21.0,3) (-28.2,3)};
  \addplot[draw=ruleGray, line width=1.3pt, forget plot] coordinates {(46.7,4) (10.0,4)};
  \addplot[draw=ruleGray, line width=1.3pt, forget plot] coordinates {(17.9,5) (-6.7,5)};
  \addplot[only marks, mark=o, mark size=2.6pt, thick, color=ink, mark options={fill=white}]
    coordinates {(33.3,1) (35.8,2) (21.0,3) (46.7,4) (17.9,5)};
  \addlegendentry{设计题}
  \addplot[only marks, mark=*, mark size=2.6pt, color=ours]
    coordinates {(-20.0,1) (-10.0,2) (-28.2,3) (10.0,4) (-6.7,5)};
  \addlegendentry{留出题}
\end{axis}
\end{tikzpicture}
\caption{从设计题到留出题的迁移。每个手段相对同批基线在通过率上的增益：空心为它所针对的设计题，实心为留出题。V1 各行
来自小样本（每边 5 到 28 次运行，见表~\ref{tab:structures}）；O2 一行比较的是多次运行结果不一致的可读题（每边 28 次试验）
与未读的 30 道题（每边两次运行）。只有验收检查者在样本外保持领先。}
\label{fig:transfer}
\end{figure}

\begin{enumerate}[leftmargin=1.6em]
  \item \textbf{同模型冗余带来的信息很少。}剩下的失败是对一句含糊的话的单一解读，而且多数运行都以同样方式解读。
    同一模型的阅读者、陪审员、评分者和检查者共享这种解读，所以它们只会确认它。这就是 \citet{huang2024cannot} 所说的
    自我纠正极限在团队层面的表现：没有新的事实，模型的第二个副本并不拥有第一个副本缺少的东西。
  \item \textbf{全局触发是用一种失败换另一种失败。}每次交付都触发的检查，修好了它所针对的题，却也返工了原本已经做对
    的题。O2 在未读题上少了 2.0，另一个全局检查少了 1.5；实时提示在它针对的 7 道题上多了 1.0，在其余 26 道上少了 1.0。
  \item \textbf{建议改变不了判断。}人设文字、规则和实时提示改变了模型下一步做什么，而不是它选择什么。阅读规则甚至让
    工作缩水，因为它让模型觉得可以跳过文字没有明说的部分。
  \item \textbf{打断缓存的 token 削减，花掉的比省下的多。}在 60k token 时移出旧工具结果，每轮省下约 10k token，却让
    未命中缓存的 token 增加 29\%、成本增加 14\%。
  \item \textbf{设计题上的胜利属于样本内。}手段是根据可读题的轨迹设计的，所以在可读题上的筛查属于样本内，会高估效果；
    这正是未读留出集所防范的自适应过拟合 \citep{dwork2015reusable}。有一个手段在筛查中 3 次通过 3 次，下一次运行时 3 次
    只通过 1 次。
\end{enumerate}

\para{我们自己的预测也是问题的一部分}每批实验之前我们都写下预测。对于根据可读题设计的手段，预测的方向常常就是错的：
O2 预测在未读题上 $+1.5$，实测 $-2.0$；实时提示预测 $+1.0$，实测 $0.0$。预先登记没有让预测变准，但让这些错误无法被
忽略。

\para{对框架构建者的启示}以下五条做法都对应上文的某项测量。
\begin{itemize}[leftmargin=1.4em]
  \item 先让提示词前缀稳定，再削减 token，并按请求测量缓存命中；打断缓存的削减花掉的比省下的多
    （第~\ref{sec:context}~节）。
  \item 看延迟时，数领队的轮次，而不是它提示词的大小（第~\ref{sec:context}~节）。
  \item 在加同一模型的智能体之前，先测一道题在多次运行之间结果不一致的频率。如果失败在每次运行中都是同一种解读，
    多几个副本只会增加成本，不会增加信息（第~\ref{sec:structure}~节和第~\ref{sec:judgment}~节）。
  \item 运行时的帮助应当由通过的运行从不出现的状态触发，例如很早就交回失败，而不是在每次交付时触发
    （第~\ref{sec:structure}~节）。
  \item 保留一个未读的留出集，让每个实验组和它的基线在同一时段运行，并在读取结果之前说明测量工具的分辨率
    （第~\ref{sec:method}~节）。
\end{itemize}

\section{有效性威胁}
\label{sec:threats}

我们的数字能分辨的差异大约是 30 道题中的两到三道，所以我们能下的结论是：我们上线或否决的任何单个手段，都没有让留出
得分变化超过这个幅度。

\begin{itemize}[leftmargin=1.4em]
  \item \textbf{噪声。}在未读的 30 道题上，两次对两次的读取，零假设标准差约为 1.4；我们在那里测量的每个单变量手段，都落在
    基线的 2 以内；只有第~\ref{sec:structure}~节的配置变化（人设与推理强度一起变）超出了这个范围。更小的真实效应，无论正负，
    都显示不出来。
  \item \textbf{在同一留出集上做了很多检验。}未读的 30 道题从未打开，但有几十个手段在它们上面读取过结果，每次读取都会
    泄露一点信息 \citep{dwork2015reusable}。一个靠运气通过的手段可能被误采用；我们采用的手段很少，主要的那一个经过了
    多次运行的确认。
  \item \textbf{抽样受额度驱动。}运行来自一个有每周额度的订阅账号。额度不足时实验暂停（10 月 26 日、11 月 8 日、
    11 月 20 日），筛查先在更便宜的模型上跑，一些实验组在已经不可能达到门槛时提前停止。
  \item \textbf{账号登录，而非 API。}几乎所有运行都使用 ChatGPT（Codex）账号登录。成本是根据服务端价格表按公开标价
    重新计算的；API 路径只通过 OpenRouter 做过一次四次尝试的冒烟测试。API 上的速率限制、延迟和缓存可能不同。
  \item \textbf{只有一台机器。}30 次尝试共用一台台式电脑，所以比在每题独占的云机器上慢；验证器在负载下有时会超时
    （由空闲重评处理）。这台电脑没有 GPU，所以 66 道题中有 3 道从未运行。
  \item \textbf{自跑的结果。}Terminal-Bench~4.1 的运行和评分都由我们自己完成，而不是由 AA 或 Terminal-Bench 维护者完成。
    适配器已经公开，以便复现（附录~\ref{app:repro}）。
  \item \textbf{两边都是我们做的。}产品、V1 任务集及其评分器都是我们的。盲做抓出了两个评分器缺陷，之后的审计也撤下或
    修复了不公平的题，但可能还有遗漏。我们的奖励投机扫描是启发式的，它标出的两次留出尝试尚未判断。
  \item \textbf{产品在变。}产品在实验期间一直在变，这正是每次比较都在同一批次内进行的原因。GPT-6.1 上的结论未必适用于
    GPT-6.2，Terminal-Bench~4.0 上的结论也未必适用于 4.1（有一道题不同）。
\end{itemize}

\section{结论与后续工作}
\label{sec:conclusion}

在一个固定的前沿模型上，一个做得好的框架能追上服务商自家的智能体，并让它更省，但还没有超过它：我们在
Terminal-Bench~4.1 上与 GPT-6.2 Sol、max 强度的 Codex 持平，输入 token 少约 30\%。

本文有三个主要结论。上下文工程的回报在成本，不在准确率：只追加的提示词把成本降低三分之一到一半，之后的手段又降低
7\% 到 20\%。结构只在工作可以拆分的地方有回报：在单个上下文装得下的任务上，一个能力强的负责者胜过所有团队，同一模型
的第二意见只会确认共同的误读，抓不住它。判断跟随模型：提示、规则和人设改变了模型下一步做什么，却没有改变它选择什么；
摆到它面前的事实和测量结果才改变了选择。

因此，要拉开差距，需要和我们已经放弃的那些手段不同类型的手段：
\begin{enumerate}[leftmargin=1.6em]
  \item \textbf{证据，而不是建议。}在决策点上，把数据的测量结果和每种待定解读的可执行后果摆到领队面前；这是唯一
    改变过选择的通道。
  \item \textbf{不共享同一解读的多样性。}同一模型的投票者在错误的解读上达成了一致。维护者榜单的第一名是一个多模型
    智能体（71.2\%），而在 Codex 中 GPT-6.2 Astra 比 GPT-6.2 Sol 高 5.5 个百分点；在解读或选择这一步加入第二个 GPT 模型，
    是显而易见的下一个实验。
  \item \textbf{更多次运行，加上有依据的选择者。}在 GPT-6.2 上，一个完美的选择者从两次运行中挑选，可以在 30 道题上多出
    约 1.6 道。选择者无论用 GPT 还是 Claude，在明显缺陷上都选得对，在需要判断的点上都选错（28 次中 5 次），所以选择者
    必须依据执行过的检查，而不是意见。
  \item \textbf{可比较的报告。}用一台有 GPU 的机器跑全部 66 道题、每题五次尝试，就能与维护者的协议一致；这是向第三方
    排行榜提交的门槛。
\end{enumerate}

\section*{数据与代码可用性}
Harbor 适配器以及复现最终运行的完整命令已经公开 \citep{zhang2026adapter}（附录~\ref{app:repro}）；已发布的服务端按
\citep{zeroth2026server} 中的文档安装。Terminal-Bench~4.1 运行的每次尝试结果（奖励、智能体用时、token 和按公开标价
计算的成本）以及可读题的轨迹由作者保存。30 道留出题的轨迹仍未读取，以便这个集合今后还能用于确认。

\section*{AI 辅助说明}
本文的实验在一个 AI 编程助手（Claude Code，模型为 Claude Opus 5.5）的帮助下设计、运行和分析，论文初稿也由它在作者
指导下撰写。作者对内容负责。

\section*{利益冲突}
作者创办了 The Zeroth，该公司开发 The One。


\bibliographystyle{plainnat}
\bibliography{refs}

\clearpage
\appendix
\section{实验索引}
\label{app:index}

本文提到的全部实验，按时间从新到旧排列。留出结果均为总和；没有点名任何留出题。“可读”和“未读”指第~\ref{sec:method}~节中
V8 的划分。

{\footnotesize
\setlength{\tabcolsep}{4pt}
\renewcommand{\arraystretch}{1.2}
\begin{longtable}{@{}>{\raggedright\arraybackslash}p{14mm} >{\raggedright\arraybackslash}p{9mm}
  >{\raggedright\arraybackslash}p{43mm} >{\raggedright\arraybackslash}p{49mm}
  >{\raggedright\arraybackslash}p{22mm}@{}}
\toprule
日期（2026） & 项目 & 实验 & 结果 & 结论 \\
\midrule
\endfirsthead
\toprule
日期（2026） & 项目 & 实验 & 结果 & 结论 \\
\midrule
\endhead
\midrule
\multicolumn{5}{r@{}}{\textit{续下页}} \\
\endfoot
\bottomrule
\endlastfoot
11/20 & 最终 & Terminal-Bench~4.1，团队 7.0，GPT-6.2 Sol，max 强度，每题 3 次 & 63 道上 60.9\%，66 道上 58.1\%；每次尝试 1.35 美元 &
  已报告 \\
11/19 & V8 & CI：在第一次写文件之后加一段提示，只给领队 & 可读 21.0 对 21.0；成本 1.00 倍 & 关闭 \\
11/18--19 & V8 & O2：首次交付时检查（GPT-6.2） & 结果不一致的可读题 $+5$；未读 $-2.0$（$p = 0.28$）；成本 1.88 倍 & 关闭 \\
11/18 & V8 & O1：为 GPT-6.2 写的一段人设 & 可读 22.0 对 20.5，一次运行；调用次数不变 & 未上线 \\
11/18 & 测试台 & 记忆教训，29 个序列 & 主助手 8/8；团队读了前一任务目录时通过（11/11 对 0/5） & 不重写记忆 \\
11/17--18 & 测试台 & 验收习惯 & 留出：团队违反习惯 14/96，主助手 0/48；误当规则 0/108 和 3/81 & 已上线 \\
11/17--18 & V8 & GPT-6.2 Sol 对 GPT-6.1 Sol & 未读 18.0 对 20.5（$p = 0.15$）；调用 31.4 对 46.9 & 基准配置保留 GPT-6.2 \\
11/15 & V8 & G：由领队启动的独立评分者 & 目标 3/10 对 5/10；10 个“全部通过”结论中 6 个落在失败的试验上 & 关闭 \\
11/14 & V8 & UL-1：按 Codex 协议主动委派 & 24 次目标中 2 次；丢两道守卫题；成本 2 到 4 倍 & 关闭；被动形式上线 \\
11/14 & V8 & CXR：动手前的阅读者 & 未读合计 $-0.8$（$\pm 0.7$） & 关闭 \\
11/13--14 & V8 & max 强度的单领队（团队 7.0）对中等强度的旧领队 & 未读 $+3.17$（标准误 1.39）；成本 2.0 倍 & 采用 \\
11/13 & V8 & E：Claude Opus 在 GPT 领队旁做阅读者或选择者 & 阅读者 0/8 且有害；选择者 38/69 & 关闭 \\
11/13 & V8 & DJ：在两份已完成构建之间选择 & 明显缺陷上选对，需要判断的点上选错；陪审团对 12、错 11、弃权 13 & 关闭 \\
11/6--8 & V8 & 关键节点重放；“先探查”文字 & 一个决策从 0 到 24/24；完整运行 $-0.04$（$\pm 0.85$） & 保留为方法 \\
11/4 & V8 & L2b，第一次留出确认 & 17/38 对 20/38 & 否决；打开 8 道留出题 \\
11 月 & V8 & GO：回应很早交回的失败 & 一道题从 0/14 到 6/6 & 已上线 \\
11 月 & V8 & HD：首次交付时检查（GPT-6.1） & 未读 $+1.54$（$\pm 1.31$）；成本 2.09 倍 & 未上线 \\
10/31 & V8 & E4，第一次 Terminal-Bench~4.0 基线（Mac） & 29/63；每次试验 1.48 美元、23.5 分钟 & 基线 \\
10/29--31 & 问答 & SWE-Atlas-QnA 全集；问题检查者；高强度 & AA 评分细则满足 70/120；检查者 $+5$，门槛 $+6$；高强度 $-2$ &
  未保留 \\
10/25--26 & 测试台 & 决策模型测试台 & 有技能时 24/24 对 16/24（$p = 0.004$） & 已上线 \\
10/25 & 测试台 & 主助手的跟进 & 评分细则 0.49 到 0.88；跟进 18/21 & 已上线 \\
10/22--24 & V1 & 进化实验室，8 轮 & 没有上线任何东西；失败来自我们自己的任务摘要 & 关闭 \\
10/19--20 & V1 & 团队 6.4 到 6.6：组合测试者、访问器规则、限定词范围 & 扩大题集后持平（16/28 对 16/28；23/34 对 24/34；23/34 对
  25/34） & 已上线 \\
10/19 & V1 & 跨模型筛查和 Claude 定夺者 & Opus 在一道题上 4/4 对 0/4；Sonnet 6/10 对 2/10 & 未上线 \\
10/19 & V1 & GPT 定夺者节点 & 合计 20/40 对 21/40 & 否决 \\
10/18 & V1 & 团队 6.3：检查者交回时的定夺步骤 & 目标 17/48 对 10/48 & 已上线 \\
10/16 & V1 & 上下文批次 b96 到 b102、多任务共用会话、回忆测试 & 采用读取指引和简短 schema；放弃精简结果和会话复用 &
  第~\ref{sec:context}~节 \\
10/15 & V1 & 精简任务板、读取工具、快速档、移出旧结果 & 都无影响；移出旧结果成本 1.14 倍 & 保留分页 \\
10/14 & V1 & 时间花在哪里 & 领队的模型调用约占活跃时间的 90\% & 提速工作针对领队轮次 \\
10/11--12 & V1 & 团队 6.0 到 6.2：验收、症状和惯例检查者 & 留出持平（5/9 对 5/9；4/8 对 5/10；6/7） & 已上线 \\
10/8--10 & V1 & 结构 1 到 10：规则、集成、阅读者、清单、架构师 & 设计题胜出，留出题落后（第~\ref{sec:discussion}~节） &
  采用检查者关口 \\
10/5 & V1 & 推理强度：中等对高 & 28/30 对 29/30，成本低 38\%；多缺陷 0/6 对 3/6 & V1 保留中等强度 \\
10/1--6 & V1 & 困难公开题和代码仓库题 & 最难的 8 道上单智能体 90.4\% 对团队 96.2\% & 扩充任务集 \\
9/30 & V1 & 只追加的纪元对旧提示词 & 成本低 33\% 到 52\%；缓存 89\%--94\% 对 76\%--81\% & 默认开启 \\
9/29 & V1 & 更大任务上的委派 & 36/36 解出；委派运行快 25\%，提示词 token 2.46 倍 & 缓存修复；未命中 $-39$\% \\
9/28 & V1 & 基线；ChatGPT 后端的缓存请求头 & 6/6 交付；缓存从 0 到 97.8\%--99.6\% & 请求头上线 \\
\end{longtable}
}

\section{复现 Terminal-Bench 4.1 运行}
\label{app:repro}

数据集为 AA 摘要版本的 Terminal-Bench~4.1，用 Harbor 0.26.0 下载：
\begin{center}
\code{terminal-bench/terminal-bench@sha256:39d9f44b\ldots e27732}
\end{center}
有 OpenAI API 密钥时，公开的适配器 \citep{zhang2026adapter} 这样运行它：

\begin{small}
\begin{verbatim}
git clone https://github.com/the-zeroth/the-one-harbor
export OPENAI_API_KEY=...
PYTHONPATH=the-one-harbor/scripts/harbor \
  uvx --python 3.12 --from harbor==0.26.0 harbor run \
  -p <dataset folder> -n 30 -k 3 \
  -a the_one_server_agent:TheOneServer -m openai/gpt-6.2-sol --effort max \
  --ak version=0.5.2 --ak run_timeout=465m
\end{verbatim}
\end{small}

\code{--ak version} 固定服务端版本，\code{--ak run\_timeout} 设定团队最长可以工作多久；服务端自己的默认值是 3 小时，
会把长任务提前截断。我们的运行有三点不同：使用产品的 ChatGPT（Codex）账号登录，通过凭据文件而不是 API 密钥；分成两个
作业（可读题和留出题）各同时运行 15 次；不含 3 道 GPU 题。Harbor 0.24 没有 \code{--effort} 选项，请改用
\code{--ak reasoning\_effort=max}。


\end{document}
